\chapter{Mesurer la latence}\label{ch:ll:measuring-latency}

Une équipe a modifié son service, relancé son test de charge et constaté une amélioration du 99\up{e}~centile. Le changement avait en fait allongé les
blocages du service, et le test avait appris à mieux ne pas les voir~: il n'envoyait chaque requête qu'après la réponse précédente,
si bien que pendant chaque blocage il cessait simplement de demander, et les requêtes qu'il aurait envoyées, celles qui auraient attendu,
n'étaient jamais mesurées. Cet effet porte un nom, l'\omterm{def:ll:measuring-latency:co}{omission coordonnée}, ainsi qu'une correction, tous deux issus d'un exposé devenu
incontournable pour quiconque publie un centile (Tene, 2015). La mesure est l'endroit où la plupart des travaux sur la latence échouent~:
l'horloge coûte autant que le code, le benchmark observe une machine chaude, l'histogramme arrondit la queue de distribution, le testeur
se coordonne avec le système qu'il teste. Ce chapitre construit l'outillage de mesure du livre et les règles
d'usage associées.

\section{Horloges}

Le chapitre~2 a introduit le \omterm{def:ll:cpu-microarchitecture:tsc}{compteur d'horodatage}, un registre qui compte à cadence constante.

\begin{definition}[Compteur d'horodatage invariant]\label{def:ll:measuring-latency:invtsc}
Un \emph{compteur d'horodatage invariant}\index{compteur d'horodatage invariant} tourne à cadence constante dans tous les états de puissance et
de fréquence du processeur, de sorte que les différences de ses lectures sont proportionnelles au temps écoulé sur chaque cœur.
Le processeur signale la propriété par un bit de CPUID~; Linux l'affiche sous forme des indicateurs \texttt{constant\_tsc} et
\texttt{nonstop\_tsc}.
\end{definition}

Avec un compteur invariant, une mesure de latence se réduit à deux lectures et une soustraction~; la conversion en nanosecondes est
une multiplication unique, par une cadence calibrée une fois pour toutes sur l'horloge monotone du système d'exploitation. Les autres solutions passent par
l'interface d'horloge du noyau, \texttt{clock\_gettime}, que Linux sert en espace utilisateur (le vDSO) lorsque sa source d'horloge
est elle-même le TSC, et sinon par un appel système. La \Cref{fig:ll:measuring-latency:clocks} mesure le coût de chaque lecture sur
ce portable~: environ \qty{6.5}{\nano\second} pour \texttt{rdtsc}, quelque \qty{11}{\nano\second} pour \texttt{rdtscp}, qui attend
les instructions précédentes, et de 14 à \qty{16}{\nano\second} pour chaque horloge du noyau. La calibration demande de la prudence quant à l'horloge du noyau qu'elle retient. \texttt{CLOCK\_MONOTONIC}, derrière \texttt{std::chrono::steady\_clock},
est discipliné par la synchronisation temporelle, qui l'accélère ou le ralentit pour converger vers l'heure exacte~;
\texttt{CLOCK\_MONOTONIC\_RAW} ne l'est pas. Dix calibrations de 50 millisecondes sur l'horloge brute concordent à six chiffres près,
soit 2.9952 ticks par nanoseconde, la fréquence nominale du processeur de 2.995~GHz~; sur l'horloge disciplinée, les mêmes intervalles donnent une
cadence d'environ 2~\% supérieure le jour de la mesure, et l'écart varie à chaque ajustement de la synchronisation. Le banc d'essai calibre
donc sur l'horloge brute.

\begin{definition}[Surcoût de mesure]\label{def:ll:measuring-latency:overhead}
Le \emph{surcoût de mesure}\index{surcoût de mesure} d'une région instrumentée est le temps que la mesure elle-même
ajoute~: les lectures d'horloge, les barrières qui les encadrent et l'enregistrement du résultat. Il est inclus dans chaque valeur relevée et doit
rester petit devant la grandeur mesurée, ou être retranché.
\end{definition}

Une étape de \qty{40}{\nano\second} mesurée avec une horloge de \qty{20}{\nano\second} est pour moitié de la mesure. Les chiffres par étape du chapitre~1
comportaient environ \qty{15}{\nano\second} de surcoût chacun~; le chapitre~26 n'horodate que les frontières d'étapes de l'ensemble du chemin, avec le
compteur, et enregistre les ticks bruts pour une conversion ultérieure.

\begin{omfigure}
\begin{tikzpicture}
\begin{axis}[width=10cm,height=5.4cm,xbar,bar width=9pt,xmin=0,enlarge x limits={upper,value=0.18},
  symbolic y coords={gettimeofday,CLOCK_REALTIME,CLOCK_MONOTONIC_RAW,CLOCK_MONOTONIC,steady_clock,rdtscp,rdtsc},
  ytick=data,yticklabels={rdtsc,rdtscp,steady\_clock,CLOCK\_MONOTONIC,CLOCK\_MONOTONIC\_RAW,CLOCK\_REALTIME,gettimeofday},
  enlarge y limits=0.12,y tick label style={font=\ttfamily\footnotesize},
  xlabel={ns par appel},tick label style={font=\small},label style={font=\small},
  grid=major,grid style={gray!25},nodes near coords,nodes near coords style={font=\footnotesize},
  every node near coord/.append style={/pgf/number format/fixed,/pgf/number format/precision=1},table/col sep=comma]
\addplot[fill=omDef!45,draw=omDef] table[y=clock,x=ns_per_call]{figdata/low-latency/05-measuring-latency/measured_clocks.csv};
\end{axis}
\end{tikzpicture}

\omcaption{Coût d'une lecture de chaque horloge accessible à un programme Linux (meilleur de sept lots d'un million d'appels, épinglé
sur un seul CPU, source d'horloge \texttt{tsc}, de sorte que chaque horloge du noyau est servie en espace utilisateur). Mesuré sur un portable (Intel Core
Ultra 7 155H) sous WSL2, sans cœurs isolés. Données~: \texttt{bench\_clocks.py}.}\label{fig:ll:measuring-latency:clocks}
\end{omfigure}

\section{Horodatage au niveau du câble}

L'horloge d'un programme ne peut dire quand le programme a vu un paquet, et non quand le paquet est arrivé.

\begin{definition}[Horodatage matériel]\label{def:ll:measuring-latency:hwts}
Un \emph{horodatage matériel}\index{horodatage matériel} est écrit par une carte réseau, un commutateur ou un dispositif de capture lorsque
les bits d'un paquet franchissent son port, à partir d'une horloge propre à ce dispositif, et il est transmis avec le paquet (par la carte) ou enregistré avec
une copie de celui-ci (par un dispositif de capture branché sur un tap ou un port miroir).
\end{definition}

L'horodatage de réception du livre~1, chapitre~28, se prend au mieux au niveau matériel~: le temps écoulé entre un \omterm{def:ll:measuring-latency:hwts}{horodatage matériel}
et la première lecture du paquet par le programme est le chemin de réception (carte, bus, noyau ou bibliothèque de contournement) que ne voit pas
un horodatage logiciel. La \omterm{def:ll:where-latency-comes-from:t2t}{latence wire-to-wire} (chapitre~1) exige deux \omterm{def:ll:measuring-latency:hwts}{horodatages matériels}, à l'entrée et à la sortie, issus d'horloges qui concordent~:
soit le même dispositif (un appareil de capture observant les deux sens), soit des dispositifs synchronisés sur une référence commune
(le protocole de temps de précision, livre~14, chapitre~4).

\section{Histogrammes et centiles}

\begin{definition}[Histogramme de latence]\label{def:ll:measuring-latency:hist}
Un \emph{histogramme de latence}\index{histogramme de latence} compte les latences enregistrées dans des classes. Dans un histogramme log-linéaire, les
classes sont exactes pour les petites valeurs et croissent en proportion de la valeur au-delà, de sorte que la largeur de chaque classe reste au plus une
fraction fixe des valeurs qu'elle contient~; les centiles sont alors connus à cette précision relative sur tout l'intervalle, depuis
une table de taille fixe et sans stocker d'échantillons.
\end{definition}

Conserver chaque échantillon d'une journée de messages est possible hors ligne, mais non sur le chemin critique (hot path)~; un histogramme de taille fixe, lui, convient.
L'histogramme du livre, \texttt{firm.lathist}, conserve 256 classes exactes puis, pour chaque puissance de deux, 128 classes~: chaque classe
a une largeur d'au plus $2^{-7}$ de ses valeurs, soit moins de 0.8~\%, depuis une nanoseconde jusqu'à la plus grande valeur sur 64~bits, dans 7\,424
compteurs.

\begin{proposition}[Précision relative des classes log-linéaires]\label{prop:ll:measuring-latency:precision}
Avec $S$ bits de sous-classes, une valeur $v \ge 2^S$ tombe dans une classe de largeur $2^{b(v)-S}$, où $b(v)$ est sa longueur en bits, et
la largeur est au plus $2^{-(S-1)}v$.
\end{proposition}

\begin{proof}
Écrivons $v = m \cdot 2^{s} + r$ avec $s = b(v) - S$, $m = \lfloor v/2^s \rfloor$ et $0 \le r < 2^s$. La classe contient les valeurs de même $m$ et de même $s$,
un intervalle de largeur $2^s$. Comme $v$ comporte $b(v)$ bits, $v \ge 2^{b(v)-1} = 2^{S-1}2^s$,
donc $2^s \le 2^{-(S-1)}v$.
\end{proof}

Les centiles se lisent dans les effectifs par le rang le plus proche~: le $p$-quantile est la borne supérieure de la classe où
l'effectif cumulé atteint pour la première fois $\lceil p\,n \rceil$ (plafonné par la plus grande valeur enregistrée). Des histogrammes issus de plusieurs
fils, machines ou journées se fusionnent par addition des effectifs, ce que les quantiles d'échantillons stockés ne permettent pas sans ces échantillons.

\section{Omission coordonnée}

\begin{definition}[Omission coordonnée]\label{def:ll:measuring-latency:co}
L'\emph{omission coordonnée}\index{omission coordonnée} est la perte des mesures les plus lentes lorsqu'un générateur de charge
attend la réponse à chaque requête avant d'envoyer la suivante~: pendant que le système se bloque, le générateur n'envoie rien, si bien que les
requêtes qui auraient attendu pendant le blocage ne sont jamais émises et jamais mesurées.
\end{definition}

\begin{omfigure}
\begin{tikzpicture}[font=\footnotesize,x=1cm,y=1cm]
\fill[omThm!10] (2,-2.3) rectangle (6,1.3);
\node[omThm,anchor=south] at (4,1.3) {service figé};
\draw[-Stealth] (0,-2.5) -- (10.4,-2.5) node[right]{temps (ms)};
\foreach \x in {0,1,2,3,4,5,6,7,8,9} {\draw (\x,-2.58) -- (\x,-2.42) node[below=3pt]{\x}; \draw[gray!40,dotted] (\x,-2.4) -- (\x,1.2);}
\node[anchor=east] at (-0.2,0.8) {boucle fermée};
\foreach \a/\b in {0/0.3,1/1.3,7/7.3,8/8.3,9/9.3} {\draw[omDef,line width=3pt] (\a,0.8) -- (\b,0.8);}
\draw[omDef,line width=3pt] (2,0.8) -- (6.1,0.8);
\node[omDef,anchor=south west] at (6.2,0.8) {un échantillon lent~; 3 à 6 jamais envoyés};
\node[anchor=east] at (-0.2,-0.4) {boucle ouverte};
\foreach \a/\b/\y in {0/0.3/-0.4,1/1.3/-0.4,2/6.1/-0.4,3/6.2/-0.8,4/6.3/-1.2,5/6.4/-1.6,6/6.5/-2.0,7/7.3/-0.4,8/8.3/-0.4,9/9.3/-0.4} {\draw[omMeth,line width=3pt] (\a,\y) -- (\b,\y);}
\node[omMeth,anchor=west] at (6.7,-1.6) {chaque requête planifiée mesurée};
\end{tikzpicture}

\omcaption{\omterm{def:ll:measuring-latency:co}{Omission coordonnée}. Chaque barre va de l'instant d'envoi prévu d'une requête jusqu'à sa réponse. Le testeur en boucle fermée (en haut) envoie à 2,
attend pendant le blocage, et n'enregistre qu'un seul échantillon lent~; les requêtes prévues à 3, 4, 5 et 6 ne sont jamais émises. Le
testeur en boucle ouverte (en bas) envoie selon le calendrier et mesure chaque requête à partir de son instant d'envoi prévu, si bien que toute requête qui
a attendu est comptée.}\label{fig:ll:measuring-latency:co-timeline}
\end{omfigure}

La correction, lorsque le rythme prévu est connu, consiste à ajouter les échantillons manquants~: une réponse ayant duré $v$ pour un intervalle prévu
$\Delta$ tient lieu de requêtes qui auraient attendu $v-\Delta$, $v-2\Delta$, et ainsi de suite jusqu'à $\Delta$. Elle est
exacte seulement si le service aurait répondu à chacune d'elles dès la fin du blocage~; une file qui se vide après le blocage
rend les valeurs réelles encore plus grandes. Le remède consiste à ne pas avoir besoin de la correction~: générer la charge selon un calendrier, en boucle ouverte, et
mesurer chaque requête à partir de l'instant où elle aurait dû être envoyée.

\begin{example}[Un blocage par seconde]\label{ex:ll:measuring-latency:co}
Un service répond en environ \qty{100}{\micro\second} et se fige pendant \qty{10}{\milli\second} au début de chaque seconde~;
les testeurs prévoient une requête par milliseconde pendant 100 secondes (la \cref{fig:ll:measuring-latency:spectrum}). Le
testeur en boucle fermée recueille 99\,000 échantillons et annonce un p99 de \qty{0.16}{\milli\second}~: le seul échantillon qu'il prélève à chaque blocage vaut 0.1~\%
du total. Corrigé, il annonce \qty{1.06}{\milli\second}~; le testeur en boucle ouverte, qui voit aussi la file se vider après
chaque blocage, \qty{1.87}{\milli\second}. Le service est lent pour 1~\% des requêtes que les utilisateurs envoient, et le test naïf
le dément d'un ordre de grandeur.
\end{example}

\begin{omfigure}
\begin{tikzpicture}
\begin{axis}[width=11cm,height=6cm,ymode=log,xmin=0.3,xmax=4,ymin=0.05,ymax=20,
  xtick={0.3010,1,2,3,4},xticklabels={50~\%,90~\%,99~\%,99.9~\%,99.99~\%},
  ytick={0.1,1,10},yticklabels={0.1,1,10},xlabel={centile},ylabel={latence (ms, échelle logarithmique)},
  tick label style={font=\small},label style={font=\small},grid=major,grid style={gray!25},
  legend columns=3,legend style={font=\footnotesize,at={(0.5,-0.3)},anchor=north,draw=none,fill=none,
  /tikz/every even column/.append style={column sep=8pt}},table/col sep=comma]
\addplot[omThm,very thick,mark=*,mark size=1.2pt] table[x=nines,y=closed_ms]{figdata/low-latency/05-measuring-latency/spectrum.csv};
\addplot[omDef,thick,dashed,mark=square*,mark size=1.2pt] table[x=nines,y=corrected_ms]{figdata/low-latency/05-measuring-latency/spectrum.csv};
\addplot[omMeth,thick,mark=triangle*,mark size=1.4pt] table[x=nines,y=open_ms]{figdata/low-latency/05-measuring-latency/spectrum.csv};
\legend{boucle fermée (brut),boucle fermée (corrigée),boucle ouverte}
\end{axis}
\end{tikzpicture}

\omcaption{Latence par centile d'un service qui se fige pendant \qty{10}{\milli\second} chaque seconde, mesurée par trois
testeurs prévoyant une requête par milliseconde. Du 99\up{e} au 99.8\up{e}~centile, le testeur en boucle fermée
se trompe d'un facteur de dix à cinquante. Données~: \texttt{ll\_measure.spectrum}, simulation déterministe de 100
secondes.}\label{fig:ll:measuring-latency:spectrum}
\end{omfigure}

Dans un système de trading, le «~testeur~» est souvent le marché. Une stratégie qui mesure son propre temps de réaction uniquement sur les
messages qu'elle parvient à traiter souffre du même biais~: pendant un blocage, les messages en attente derrière elle sont traités en retard,
et une mesure qui horodate chaque message au moment où le programme le touche pour la première fois, plutôt qu'à l'arrivée du paquet,
n'enregistre aucun de ce retard. L'horodatage de réception relève du support physique (un \omterm{def:ll:measuring-latency:hwts}{horodatage matériel}) ou, du moins, au
moment où le paquet a quitté le noyau.

\section{Profileurs et compteurs matériels}

\begin{definition}[Profileur par échantillonnage, compteur de performance matériel]\label{def:ll:measuring-latency:prof}
Un \emph{profileur par échantillonnage}\index{profileur par échantillonnage} interrompt un programme à intervalles réguliers (de temps ou
d'événements matériels) et enregistre l'endroit où il se trouvait, en estimant la part du temps passée dans chaque fonction. Un \emph{compteur de performance
matériel}\index{compteur de performance matériel} est un registre du processeur qui compte des événements (cycles, instructions
retirées, \omterm{def:ll:memory-hierarchy-and-caches:miss}{défauts de cache}, \omterm{def:ll:cpu-microarchitecture:branch}{erreurs de prédiction de branchement}) pour le code exécuté sur un cœur~; sous Linux, elles se lisent à l'aide de
\texttt{perf\_event\_open} et l'outil \texttt{perf}.
\end{definition}

Les deux répondent à la question «~où passe le temps en moyenne~», qui est la mauvaise question pour une queue de distribution. Un \omterm{def:ll:measuring-latency:prof}{profileur par échantillonnage} ne voit
rien d'un événement qui survient une fois par minute~; les compteurs ne donnent que des totaux, sauf lus autour d'une région. Leur usage dans le travail sur la
latence relève du diagnostic~: compter les \omterm{def:ll:memory-hierarchy-and-caches:miss}{défauts de cache} et les erreurs de prédiction dans la région que les horodatages signalent comme lente, sur une
reproduction du cas lent. L'accès est souvent restreint~: le paramètre du noyau \texttt{perf\_event\_paranoid}, égal à 2 par défaut,
n'autorise que les mesures en espace utilisateur, et les machines virtuelles, y compris l'ordinateur portable de ce livre, n'exposent en général
pas du tout les compteurs. Les mesures de ce livre n'utilisent donc que l'horloge et ce qui peut être compté en logiciel.

\begin{method}[Rapporter une mesure de latence]\label{met:ll:measuring-latency:report}
\begin{enumerate}
\item Nommer les deux instants et préciser qui a horodaté chacun (la \cref{def:ll:where-latency-comes-from:t2t} au chapitre~1).
\item Piloter le système selon un calendrier (en boucle ouverte) à un rythme indiqué, ou corriger l'\omterm{def:ll:measuring-latency:co}{omission coordonnée} et le dire.
\item Enregistrer chaque mesure dans un histogramme à erreur relative bornée~; annoncer le p50, le p99, le p99.9 et le maximum, avec le
nombre de mesures.
\item Indiquer la machine, le placement, le compilateur et ses options, l'horloge et son surcoût, l'échauffement, et les autres processus en cours d'exécution.
\item Répéter la mesure, et annoncer la dispersion des centiles entre les exécutions avant de comparer deux versions (chapitre~25).
\end{enumerate}
\end{method}

\section{Tutoriel~: un histogramme et un générateur de charge}\label{tut:ll:measuring-latency:tut}

\begin{tutorial}
\textbf{Objectif.} Construire l'histogramme log-linéaire dans trois langages sur un même jeu de test, mesurer les horloges et reproduire
l'\omterm{def:ll:measuring-latency:co}{omission coordonnée}. \textbf{État final~:} les \cref{fig:ll:measuring-latency:clocks,fig:ll:measuring-latency:spectrum}.
\begin{tutsteps}
\item \textbf{Indexer une valeur.} En dessous de $2^S$, la valeur est son propre indice~; au-dessus, ses $S$ bits de poids fort et sa longueur en bits choisissent la
case.
\omcode{code/firm/lathist/cpp/firm_lathist.hpp}{12}{18}{La case d'une valeur~: une longueur en bits et un décalage.}\label{lst:ll:measuring-latency:index}
\item \textbf{Corriger l'omission.} L'enregistrement corrigé ajoute les échantillons que le blocage a masqués.
\omcode{code/firm/lathist/cpp/firm_lathist.hpp}{39}{44}{Enregistrement avec l'intervalle attendu.}\label{lst:ll:measuring-latency:corrected}
\item \textbf{Vérifier les jumeaux.} \texttt{make\_lathist\_fixture.py} écrit 20\,000 valeurs allant d'une nanoseconde à dix secondes,
leurs cases non nulles et sept quantiles~; les tests Python, C++20 et Rust reproduisent les deux fichiers à l'identique.
\item \textbf{Simuler les testeurs.} Le testeur en boucle fermée envoie dès l'arrivée de la réponse précédente~; celui en boucle ouverte, à la
milliseconde. \texttt{fig\_measure.py} écrit le spectre.
\omcode{code/low-latency/05-measuring-latency/python/ll_measure.py}{45}{53}{Un testeur en boucle fermée, source de l'omission coordonnée.}\label{lst:ll:measuring-latency:closed}
\item \textbf{Mesurer les horloges} avec \texttt{python bench\_clocks.py}.
\end{tutsteps}
\textbf{Pour aller plus loin.} Faire durer le blocage \qty{50}{\milli\second} une fois toutes les cinq secondes et comparer les trois
p99~; ajouter un \texttt{record\_corrected} au banc d'essai du chapitre~1, avec l'intervalle d'un flux réel, et voir si quelque chose
change.
\end{tutorial}

\section{Construction~: l'histogramme de latence}\label{bld:ll:measuring-latency:lathist}

\begin{build}
\bfield{Rôle} Le seul moyen dont dispose la firme pour enregistrer une distribution de latence, sur le chemin critique (hot path) comme en dehors de celui-ci, dans tous les langages, afin que
les histogrammes issus du feed handler, du moteur de stratégie et de la passerelle puissent être fusionnés et comparés.
\bfield{Interface} Python \texttt{LatHist(sub\_bits)} avec \texttt{record(v, n)},
\texttt{record\_corrected(v, interval)}, \texttt{merge}, \texttt{quantile(p)}, \texttt{mean}, \texttt{nonzero}~; \texttt{index\_of}, \texttt{bucket\_range}.
C++20 \texttt{firm::lathist::LatHist} et Rust \texttt{firm\_lathist::LatHist}~: les mêmes, avec les compteurs sous forme de
vecteur de taille fixe.
\bfield{Règles} Les valeurs sont des entiers non négatifs (nanosecondes, ou ticks convertis une seule fois)~; la largeur d'une case vaut au plus
$2^{-(S-1)}$ de ses valeurs~; les quantiles se calculent par rang le plus proche à partir du bord supérieur de la case, plafonnés par le maximum~; l'enregistrement n'alloue jamais
de mémoire~; les trois implémentations concordent exactement sur la fixture.
\bfield{Tests d'acceptation} \texttt{code/firm/lathist/}~: les cases qui recouvrent la droite, la borne d'erreur relative, les quantiles et
la fusion, la correction d'omission, et la fixture commune en Python, C++20 et Rust.
\bfield{Pour aller plus loin} Sérialiser un histogramme de façon compacte pour la journalisation (longueurs de séries pour les cases nulles)~; un histogramme d'intervalle qui
tourne chaque seconde sans interrompre l'enregistreur.
\end{build}

\begin{omsources}
\item G.~Tene, ``How NOT to measure latency'', QCon San Francisco 2015.
\item HdrHistogram, project documentation (``Corrected vs.\ raw value recording calls'').
\item Intel, \emph{Intel 64 and IA-32 Architectures Software Developer's Manual}, vol.~3, ``Invariant TSC''.
\item Linux man pages \texttt{clock\_gettime(2)}, \texttt{vdso(7)}, \texttt{perf\_event\_open(2)}.
\end{omsources}

\section{Exercices}

\begin{exercise}[$\star$]\label{exo:ll:measuring-latency:1}
Une région est chronométrée à 3\,020 ticks sur un compteur calibré à 3.02 ticks par nanoseconde, avec un surcoût de 45 ticks. Quelle
est sa latence~?
\end{exercise}

\begin{exercise}[$\star$]\label{exo:ll:measuring-latency:2}
Quelle case de \texttt{firm.lathist} (avec $S = 8$) contient \qty{1\,000}{\nano\second}, et quelle plage de valeurs
couvre-t-elle~?
\end{exercise}

\begin{exercise}[$\star$]\label{exo:ll:measuring-latency:3}
Combien de compteurs l'histogramme nécessite-t-il avec $S = 10$, et quelle est sa précision relative~?
\end{exercise}

\begin{exercise}[$\star\star$]\label{exo:ll:measuring-latency:4}
Un testeur en boucle fermée qui vise une requête toutes les \qty{2}{\milli\second} enregistre une réponse de \qty{9}{\milli\second}. Quelles
valeurs la correction ajoute-t-elle~?
\end{exercise}

\begin{exercise}[$\star\star$]\label{exo:ll:measuring-latency:5}
Pourquoi le testeur en boucle ouverte de l'\cref{ex:ll:measuring-latency:co} affiche-t-il un p99 plus élevé que le
testeur en boucle fermée corrigé~?
\end{exercise}

\begin{exercise}[$\star\star$]\label{exo:ll:measuring-latency:6}
La durée de réaction d'une stratégie est mesurée à partir de l'instant où sa boucle d'événements lit chaque message. Pendant un
blocage de \qty{2}{\milli\second}, 400 messages s'accumulent dans le socket. Qu'enregistre la mesure pour eux, et que devrait-elle enregistrer~?
\end{exercise}

\begin{exercise}[$\star\star\star$]\label{exo:ll:measuring-latency:7}
\emph{Programmation.} Exécuter \texttt{ll\_measure} avec un blocage de \qty{50}{\milli\second} toutes les cinq secondes. Quels sont les trois p99
désormais, et pourquoi le chiffre en boucle fermée bouge-t-il à peine~?
\end{exercise}

\begin{exercise}[$\star\star\star$]\label{exo:ll:measuring-latency:8}
\emph{Trouver la faille.} «~Le p99.99 de notre service est de \qty{40}{\micro\second}~: nous avons chronométré un million de requêtes depuis un seul
thread client, en envoyant chaque requête dès que la réponse précédente est arrivée, et nous avons conservé les échantillons dans un histogramme de 1\,000
cases de \qty{100}{\nano\second}.~»
\end{exercise}

\section{Problème~: le banc d'essai qui mentait}

\begin{problem}[{Problème du week-end --- un blocage par seconde}]\label{pb:ll:measuring-latency:1}
Le service de l'\cref{ex:ll:measuring-latency:co} répond en environ \qty{100}{\micro\second} et se bloque pendant
\qty{10}{\milli\second} au début de chaque seconde. Trois testeurs visent une requête par milliseconde pendant 100 secondes.

\medskip\noindent\textbf{Partie I --- Comptage.}
\begin{enumerate}
\item Combien de requêtes les testeurs en boucle ouverte et en boucle fermée envoient-ils, et pourquoi ce nombre diffère-t-il~?
\item Quelle part des échantillons en boucle fermée est lente~?
\item Quelle part des requêtes envoyées par un client en boucle ouverte arrive pendant un blocage ou attend derrière un blocage~?
\item Combien d'échantillons la correction ajoute-t-elle par blocage, et au total~?
\end{enumerate}

\medskip\noindent\textbf{Partie II --- Centiles.}
\begin{enumerate}[resume]
\item Lire les trois p99 à partir de la \cref{fig:ll:measuring-latency:spectrum}.
\item Pourquoi les trois testeurs concordent-ils à la médiane et au p99.9~?
\item Pourquoi le chiffre corrigé se situe-t-il en dessous de celui en boucle ouverte~?
\item Quel chiffre un utilisateur du service ressentirait-il~?
\end{enumerate}

\medskip\noindent\textbf{Partie III --- L'histogramme.}
\begin{enumerate}[resume]
\item Quelle est la précision relative de l'histogramme utilisé, et cela importe-t-il ici~?
\item Combien de compteurs contient-il~?
\item Les trois histogrammes des testeurs pourraient-ils être fusionnés, et que signifierait l'histogramme fusionné~?
\item Que donnerait un histogramme à cases fixes de \qty{100}{\micro\second} jusqu'à \qty{10}{\milli\second}~?
\end{enumerate}

\medskip\noindent\textbf{Partie IV --- Le rapport.}
\begin{enumerate}[resume]
\item Énoncer le \emph{résultat nommé}~: le p99 rapporté par le testeur en boucle fermée, brut et corrigé, face à celui du
testeur en boucle ouverte.
\item Réécrire le rapport naïf sous la forme de la \cref{met:ll:measuring-latency:report}.
\item Pourquoi «~un blocage par seconde~» est-il réaliste pour un service doté d'un ramasse-miettes ou d'une tâche périodique~?
\item Comment détecter l'\omterm{def:ll:measuring-latency:co}{omission coordonnée} dans un rapport qu'on n'a pas produit~?
\item Quel est l'analogue, sur le marché, d'un testeur en boucle fermée~?
\item Pourquoi le surcoût de l'horloge doit-il être indiqué avec chaque chiffre~?
\item Qu'apporte un \omterm{def:ll:measuring-latency:hwts}{horodatage matériel} à la mesure d'un feed handler~?
\item En une phrase~: que dit un chiffre de latence honnête~?
\end{enumerate}
\end{problem}

\section{Questions d'entretien}

\begin{interviewq}[$\star$ \iqroles{développeur}]\label{iq:ll:measuring-latency:1}
Comment lisez-vous l'heure avec le moindre surcoût sur x86 Linux, et que doit-il être vrai pour que la lecture ait un sens
d'un cœur à l'autre~?
\end{interviewq}

\begin{interviewq}[$\star\star$ \iqroles{développeur}]\label{iq:ll:measuring-latency:2}
Qu'est-ce que l'\omterm{def:ll:measuring-latency:co}{omission coordonnée}, et comment concevriez-vous un test de charge qui l'évite~?
\end{interviewq}

\begin{interviewq}[$\star\star$ \iqroles{développeur}]\label{iq:ll:measuring-latency:3}
Vous devez enregistrer le p99.9 d'un flux d'un million de latences par seconde à mémoire constante. Comment faites-vous~?
\end{interviewq}

\begin{interviewq}[$\star\star$ \iqroles{développeur, chercheur}]\label{iq:ll:measuring-latency:4}
Pourquoi la latence moyenne est-elle un mauvais résumé, et que rapporteriez-vous à la place~?
\end{interviewq}

\begin{interviewq}[$\star\star$ \iqroles{développeur}]\label{iq:ll:measuring-latency:5}
Quand un \omterm{def:ll:measuring-latency:prof}{profileur par échantillonnage} est-il utile pour le travail sur la latence, et quand est-il inutile~?
\end{interviewq}

\begin{interviewq}[$\star\star\star$ \iqroles{développeur}]\label{iq:ll:measuring-latency:6}
Deux versions d'un moteur de trading diffèrent de 3~\% au p99 sur un seul essai chacune. De quoi avez-vous besoin avant d'y croire~?
\end{interviewq}
